\documentclass[10.5pt,fontset=fandol]{ctexart}
\usepackage[paperwidth=222.6413mm,paperheight=297mm,left=1.9cm,right=1.9cm,top=2.1cm,bottom=1.8cm,headsep=0.5cm]{geometry}
\usepackage{amsmath,amssymb,mathrsfs}
\usepackage[table]{xcolor}
\usepackage{booktabs,tabularx,xltabular,needspace,array,ragged2e,enumitem,tcolorbox,tikz,fancyhdr}
\tcbuselibrary{skins,breakable}
\usepackage[hidelinks]{hyperref}
\definecolor{deepblue}{HTML}{1B3A5C}\definecolor{navy}{HTML}{1B3A5C}\definecolor{accentred}{HTML}{B23A48}\definecolor{accentA}{HTML}{B23A48}
\definecolor{muted}{HTML}{5A6470}\definecolor{qbg}{HTML}{F4F6F8}\definecolor{linec}{HTML}{9AA3AD}
\newcolumntype{L}[1]{>{\RaggedRight\arraybackslash}p{#1}}
\newlist{qp}{enumerate}{1}\setlist[qp]{label=(\alph*),leftmargin=1.5em,itemsep=2pt}
\newlist{qn}{enumerate}{1}\setlist[qn]{label=(\arabic*),leftmargin=1.5em,itemsep=2pt}
\newlist{qparts}{enumerate}{1}\setlist[qparts]{label=(\alph*),leftmargin=1.5em,itemsep=2pt}
\newlist{qnum}{enumerate}{1}\setlist[qnum]{label=(\arabic*),leftmargin=1.5em,itemsep=2pt}
\newlist{bul}{itemize}{1}\setlist[bul]{label=\textbullet,leftmargin=1.5em,itemsep=2pt}
\newcommand{\key}[1]{\par\vspace{4pt}\noindent\textbf{\color{navy}#1}\par\vspace{2pt}}
\newcommand{\hint}[1]{\par\smallskip{\small\itshape Hint: #1}\par}
\newcommand{\sect}[1]{{\small\color{muted}#1}}
\newcommand{\dd}{\mathrm d}\newcommand{\Var}{\operatorname{Var}}\newcommand{\Cov}{\operatorname{Cov}}
\newcommand{\E}{\mathbb E}\newcommand{\ee}{\mathrm e}\newcommand{\jj}{\mathrm j}
\newtcolorbox{workspace}{height fill,enhanced,colback=white,colframe=linec!55,boxrule=0.5pt,arc=1.5pt}
\newcommand{\qtitle}[2]{\par\addvspace{1.2em}\Needspace{8\baselineskip}\begin{tcolorbox}[colback=navy,colframe=navy,arc=1.5pt,boxrule=0pt]\color{white}\large\bfseries #1\hfill\small #2\end{tcolorbox}}
\pagestyle{fancy}\fancyhf{}\fancyhead[L]{\footnotesize\color{navy}SSA -- Chapter 1 -- Solutions}\fancyhead[R]{\footnotesize\thepage}
\fancyfoot[C]{\scriptsize Eight-paper edition -- 2026-10-13}
\begin{document}

\begin{titlepage}\centering\vspace*{1.5cm}
{\Huge\bfseries 随机信号分析}\par\vspace{1cm}
{\LARGE\bfseries 第1章复习习题集中文详解}\par\vspace{1cm}
{\large 八套真题版 · 30题 · 10天 · 每天3题}\par\vspace{1cm}
\begin{tcolorbox}[colback=qbg,colframe=navy]
本章八套独占核心均值为27.500分/套；最近四套为33.75分/套。解答按配套英文题集编号，每题保留来源及教材小节。原题条件不足时明确写出可确定的结论与额外假设；奇异高斯分布不强套普通PDF。
\end{tcolorbox}\vfill 2026-10-13\end{titlepage}
\section*{题号与来源索引}
{\small\renewcommand{\arraystretch}{1.15}\begin{xltabular}{\textwidth}{@{}r r L{7cm} X@{}}\toprule 题 & 天 & 题源 & 教材小节\\\midrule\endfirsthead\toprule 题 & 天 & 题源 & 教材小节\\\midrule\endhead\midrule\multicolumn{4}{r}{\small 索引续下页}\\\endfoot\bottomrule\endlastfoot
1 & 1 & Formula recall \textemdash\ \S1.4.1 and \S1.4.2 & 1.4.1 / 1.4.2\\
2 & 1 & Textbook 1.21 / Required & 1.4.2\\
3 & 1 & 2023--24 Q3(a)--(c); same density as 2018--19 Q2 & 1.2.1 / 1.4.1 / 1.4.2\\
4 & 2 & Final examination 2023--24, Q3(d) & 1.4.2 / 1.4.3\\
5 & 2 & Textbook Exercise 1.22 & 1.4.3\\
6 & 2 & Final examination 2024--25, Q1(b) & 1.4.3 / 1.2.3\\
7 & 3 & Textbook 1.14 / Required & 1.2.1 / 1.4.2\\
8 & 3 & Final examination 2025--26, Q1(a)--(c) & 1.2.1 / 1.4.2\\
9 & 3 & Final examination 2025--26, Q1(d) & 1.2.1 / 1.2.3\\
10 & 4 & Formula recall \textemdash\ \S1.5.1 & 1.5.1\\
11 & 4 & Final examination 2022--23, Q1(b) & 1.5.1\\
12 & 4 & Textbook 1.29 / Optional & 1.5.1\\
13 & 5 & Formula recall \textperiodcentered\ \S1.6 & 1.6\\
14 & 5 & Textbook Exercise 1.17 & 1.3.1\\
15 & 5 & Adapted from 2023--24 Q2(c) & 1.6 / 1.4.2\\
16 & 6 & Textbook 1.16 / Required & 1.3.2\\
17 & 6 & Formula recall \textperiodcentered\ \S1.2.1 and \S1.6 & 1.2.1 / 1.6\\
18 & 6 & Final examination 2025--26, Q4(b)--(c) & 1.6 / 1.4.2\\
19 & 7 & Textbook 1.6 / Required & 1.1.1\\
20 & 7 & Textbook 1.24 / Required & 1.4.4\\
21 & 7 & Textbook 1.33 / Required & 1.5.2\\
22 & 8 & 2018--19 Q1 (201812) & 1.3.1 / 1.6\\
23 & 8 & Adapted from 2019--20 Q1 (201912) & 1.3.1 / 1.5.1 / 1.6\\
24 & 8 & 2019--20 Q2 (201912) & 1.2.1\\
25 & 9 & Adapted from 2020--21 Q4(b,c) (202104) & 1.6 / 2.4.3\\
26 & 9 & 2021--22 Q1(c) (202112) & 1.5.2 / 1.4.2\\
27 & 9 & Textbook 1.9 / Optional & 1.1.2\\
28 & 10 & Textbook 1.11 / Optional & 1.1.2\\
29 & 10 & Textbook 1.28 / Optional & 1.5.1\\
30 & 10 & Textbook 1.30 / Optional & 1.5.1\\
\end{xltabular}}\par\smallskip 习题与详解题号完全一致；Day标签用于安排，不表示题目在原卷中的顺序。\clearpage
\qtitle{第 1 题}{Day 1}
{\small\color{accentred}(Formula recall \textemdash\ \S1.4.1 and \S1.4.2)}\par\smallskip
{\small\color{muted}教材小节：1.4.1 / 1.4.2}\par\medskip
\noindent\textbf{数字特征的定义与 $D(aX{+}bY)$ 的推导}\par\smallskip\begin{tcolorbox}[breakable,colback=qbg,colframe=qbg,boxrule=0pt]设 $X$ 的概率密度为 $f(x)$，二维情形为 $f(x,y)$；离散型把积分换成求和即可。

\key{七个定义}
\begin{qp}
  \item $n$ 阶原点矩：$m_n=E[X^n]=\displaystyle\int_{-\infty}^{+\infty}x^nf(x)\dd x$
  \item $n$ 阶中心矩：$\mu_n=E\big[(X-EX)^n\big]$
  \item 均方值 $E[X^2]$；均方根值（有效值）$\sqrt{E[X^2]}$
  \item 方差 $D(X)=E\big[(X-EX)^2\big]=E[X^2]-(EX)^2$；标准差 $\sigma_X=\sqrt{D(X)}$
  \item 联合原点矩：$m_{kl}=E[X^kY^l]=\displaystyle\iint x^ky^lf(x,y)\dd x\dd y$
  \item 协方差：$\mathrm{Cov}(X,Y)=E\big[(X-EX)(Y-EY)\big]=E[XY]-EX\cdot EY$
  \item 相关系数：$\rho_{XY}=\dfrac{\mathrm{Cov}(X,Y)}{\sigma_X\sigma_Y}$，且 $|\rho_{XY}|\le1$
\end{qp}

\key{推导 $D(aX+bY)$}
记 $U=aX+bY$，则
\[
U-EU=a(X-EX)+b(Y-EY).
\]
两边取期望的平方（注意 $E[X-EX]=E[Y-EY]=0$）：
\[
D(U)=E\big[(U-EU)^2\big]
=a^2E\big[(X-EX)^2\big]+b^2E\big[(Y-EY)^2\big]+2ab\,E\big[(X-EX)(Y-EY)\big],
\]
即
\[
\boxed{\,D(aX+bY)=a^2D(X)+b^2D(Y)+2ab\,\mathrm{Cov}(X,Y)\,}
\]
当 $X,Y$ 不相关（$\mathrm{Cov}=0$）时退化为 $D(aX+bY)=a^2D(X)+b^2D(Y)$；
当 $a=1,b=\pm1$ 时即为常用的 $D(X\pm Y)=D(X)+D(Y)\pm2\mathrm{Cov}(X,Y)$。
\end{tcolorbox}
\qtitle{第 2 题}{Day 1}
{\small\color{accentred}(Textbook 1.21 / Required)}\par\smallskip
{\small\color{muted}教材小节：1.4.2}\par\medskip
\noindent\textbf{线性组合的 $E$、$D$、$\mathrm{Cov}$}\par\smallskip\begin{tcolorbox}[breakable,colback=qbg,colframe=qbg,boxrule=0pt]已知 $EX=1$，$EY=3$，$D(X)=4$，$D(Y)=16$，$\rho_{XY}=0.5$。

\key{先求协方差}
由 $\rho_{XY}=\dfrac{\mathrm{Cov}(X,Y)}{\sigma_X\sigma_Y}$ 得
\[
\mathrm{Cov}(X,Y)=\rho_{XY}\sqrt{D(X)}\sqrt{D(Y)}=0.5\times2\times4=4 .
\]

\key{期望}
\[
EU=E(3X+Y)=3EX+EY=3(1)+3=6,
\]
\[
EV=E(X-2Y)=EX-2EY=1-6=-5 .
\]

\key{方差}
由 $D(aX+bY)=a^2D(X)+b^2D(Y)+2ab\,\mathrm{Cov}(X,Y)$：
\[
D(U)=D(3X+Y)=9D(X)+D(Y)+2\cdot3\cdot1\cdot\mathrm{Cov}(X,Y)
=9(4)+16+6(4)=36+16+24=76,
\]
\[
D(V)=D(X-2Y)=D(X)+4D(Y)+2\cdot1\cdot(-2)\cdot\mathrm{Cov}(X,Y)
=4+64-16=52 .
\]

\key{协方差}
由 $\mathrm{Cov}(aX+bY,\;cX+dY)=acD(X)+ad\,\mathrm{Cov}(X,Y)+bc\,\mathrm{Cov}(Y,X)+bdD(Y)$，
取 $a=3,b=1,c=1,d=-2$：
\[
\mathrm{Cov}(U,V)=3(4)+3(-2)(4)+1(1)(4)+1(-2)(16)=12-24+4-32=-40 .
\]

\key{校验}
用 $E[UV]-EU\cdot EV$ 复核：$UV=(3X+Y)(X-2Y)=3X^2-5XY-2Y^2$，
$E[X^2]=D(X)+(EX)^2=5$，$E[Y^2]=16+9=25$，$E[XY]=\mathrm{Cov}+EX\cdot EY=7$，故
$E[UV]=15-35-50=-70$，$\mathrm{Cov}(U,V)=-70-6\times(-5)=-40$ \checkmark
\end{tcolorbox}
\qtitle{第 3 题}{Day 1}
{\small\color{accentred}(2023--24 Q3(a)--(c); same density as 2018--19 Q2)}\par\smallskip
{\small\color{muted}教材小节：1.2.1 / 1.4.1 / 1.4.2}\par\medskip
\noindent\textbf{联合密度归一化、期望与方差}\par\smallskip\begin{tcolorbox}[breakable,colback=qbg,colframe=qbg,boxrule=0pt]\[
f_{XY}(x,y)=A\sin(x+y),\qquad 0\le x\le \frac{\pi}{2},\qquad 0\le y\le\frac{\pi}{2}.
\]

\key{(a) 求 $A$}
先对 $y$ 积分：
\[
\int_0^{\pi/2}\sin(x+y)\dd y=\big[-\cos(x+y)\big]_0^{\pi/2}=\sin x+\cos x .
\]
再对 $x$ 积分：
\[
\iint f_{XY}\dd x\dd y=A\int_0^{\pi/2}(\sin x+\cos x)\dd x
=A\big[-\cos x+\sin x\big]_0^{\pi/2}=A(1+1)=2A .
\]
由归一化 $2A=1$ 得
\[
\boxed{A=\tfrac12}
\]

\key{(b) 期望}
用同一个内层积分 $\int_0^{\pi/2}\sin(x+y)\dd y=\sin x+\cos x$：
\[
E(X)=\frac12\int_0^{\pi/2}x(\sin x+\cos x)\dd x .
\]
其中
\[
\int_0^{\pi/2}x\sin x\dd x=\big[-x\cos x+\sin x\big]_0^{\pi/2}=1,
\qquad
\int_0^{\pi/2}x\cos x\dd x=\big[x\sin x+\cos x\big]_0^{\pi/2}=\frac{\pi}{2}-1 .
\]
两式相加得 $\dfrac{\pi}{2}$，故
\[
EX=\frac12\cdot\frac{\pi}{2}=\frac{\pi}{4}\approx0.7854 .
\]
由于 $f_{XY}$ 关于 $x,y$ 完全对称，同理
\[
\boxed{EX=EY=\tfrac{\pi}{4}}
\]

\key{(c) 方差}
\[
E(X^2)=\frac12\int_0^{\pi/2}x^2(\sin x+\cos x)\dd x .
\]
由 $\int x^2\sin x\dd x=-x^2\cos x+2x\sin x+2\cos x$ 得 $\int_0^{\pi/2}x^2\sin x\dd x=\pi-2$；
由 $\int x^2\cos x\dd x=x^2\sin x+2x\cos x-2\sin x$ 得 $\int_0^{\pi/2}x^2\cos x\dd x=\dfrac{\pi^2}{4}-2$。
两者相加得 $\dfrac{\pi^2}{4}+\pi-4$，故
\[
E(X^2)=\frac12\Big(\frac{\pi^2}{4}+\pi-4\Big)=\frac{\pi^2}{8}+\frac{\pi}{2}-2 .
\]
\[
\boxed{D(X)=E(X^2)-(EX)^2=\frac{\pi^2}{16}+\frac{\pi}{2}-2\approx0.1876}
\]
同理 $D(Y)=D(X)$。
\end{tcolorbox}
\qtitle{第 4 题}{Day 2}
{\small\color{accentred}(Final examination 2023--24, Q3(d))}\par\smallskip
{\small\color{muted}教材小节：1.4.2 / 1.4.3}\par\medskip
\noindent\textbf{联合矩、相关系数与独立/不相关/正交的判定}\par\smallskip\begin{tcolorbox}[breakable,colback=qbg,colframe=qbg,boxrule=0pt]沿用第 3 题的 $f_{XY}(x,y)=\dfrac12\sin(x+y)$，$0\le x,y\le\dfrac{\pi}{2}$。

\key{(a) 联合矩 $m_{XY}=E(XY)$}
利用 $\sin(x+y)=\sin x\cos y+\cos x\sin y$，被积函数可分离：
\[\begin{aligned}
E(XY)&=\frac12\iint xy\sin(x+y)\,dx\,dy\\
&=\frac12(I_sI_c+I_cI_s),\\
I_s&=\int_0^{\pi/2}x\sin x\,dx,\qquad
I_c=\int_0^{\pi/2}x\cos x\,dx.
\end{aligned}\]
代入第 3 题算过的 $\int_0^{\pi/2}x\sin x\dd x=1$、$\int_0^{\pi/2}x\cos x\dd x=\dfrac{\pi}{2}-1$：
\[
E(XY)=\frac12\Big[1\cdot\Big(\frac{\pi}{2}-1\Big)+\Big(\frac{\pi}{2}-1\Big)\cdot1\Big]
=\frac{\pi}{2}-1\approx0.5708 .
\]

\key{(b) 相关系数 $\rho_{XY}$}
\[
\mathrm{Cov}(X,Y)=E(XY)-EX\cdot EY=\Big(\frac{\pi}{2}-1\Big)-\frac{\pi^2}{16}\approx-0.0461,
\]
由第 3 题 $\sigma_X=\sigma_Y=\sqrt{0.187646}\approx0.4332$，
\[
\boxed{\rho_{XY}=\frac{\mathrm{Cov}(X,Y)}{\sigma_X\sigma_Y}
=\frac{\frac{\pi}{2}-1-\frac{\pi^2}{16}}{\frac{\pi^2}{16}+\frac{\pi}{2}-2}\approx-0.2454}
\]

\key{三个判定}
\begin{bul}
  \item \textbf{不相关？不成立。} $\rho_{XY}\approx-0.2454\neq0$，说明 $X$ 与 $Y$ \textbf{相关}。
  \item \textbf{正交？不成立。} 正交要求 $E[XY]=0$，而 $E[XY]\approx0.5708\neq0$。
        根源是 $EX=EY=\tfrac{\pi}{4}\neq0$——"不相关"并不能推出"正交"。
  \item \textbf{独立？不成立。} 由边缘密度
        $f_X(x)=\int_0^{\pi/2}\tfrac12\sin(x+y)\dd y=\tfrac12(\sin x+\cos x)$（$f_Y$ 同理），
        \[
        f_X(x)f_Y(y)=\tfrac14(\sin x+\cos x)(\sin y+\cos y)\neq\tfrac12\sin(x+y)=f_{XY}(x,y),
        \]
        故不独立。（也可用"独立 $\Rightarrow$ 不相关"的逆否命题，由 $\rho\neq0$ 直接否定独立性。）
\end{bul}
\end{tcolorbox}
\qtitle{第 5 题}{Day 2}
{\small\color{accentred}(Textbook Exercise 1.22)}\par\smallskip
{\small\color{muted}教材小节：1.4.3}\par\medskip
\noindent\textbf{正交 / 不相关 / 独立：独立但不正交的例子}\par\smallskip\begin{tcolorbox}[breakable,colback=qbg,colframe=qbg,boxrule=0pt]\[
f_{XY}(x,y)=\frac{xy}{9},\qquad 0\le x<2,\quad 0\le y<3 .
\]

\key{先验证归一化}
\[
\iint f_{XY}\dd x\dd y=\frac19\Big(\int_0^2x\dd x\Big)\Big(\int_0^3y\dd y\Big)=\frac19\cdot2\cdot\frac92=1 \quad\text{\checkmark}
\]

\key{边缘密度}
\[
f_X(x)=\int_0^3\frac{xy}{9}\dd y=\frac{x}{9}\cdot\frac92=\frac{x}{2}\ (0<x<2),
\qquad
f_Y(y)=\int_0^2\frac{xy}{9}\dd x=\frac{y}{9}\cdot2=\frac{2y}{9}\ (0<y<3).
\]

\key{独立？\ 是}
\[
f_X(x)f_Y(y)=\frac{x}{2}\cdot\frac{2y}{9}=\frac{xy}{9}=f_{XY}(x,y),
\]
联合密度可分解为边缘密度之积，故 $X$ 与 $Y$ \textbf{相互独立}。

\key{不相关？\ 是}
\[
EX=\int_0^2x\cdot\frac{x}{2}\dd x=\frac43,\qquad
EY=\int_0^3y\cdot\frac{2y}{9}\dd y=2,
\]
\[
E(XY)=\frac19\Big(\int_0^2x^2\dd x\Big)\Big(\int_0^3y^2\dd y\Big)=\frac19\cdot\frac83\cdot9=\frac83,
\]
\[
\mathrm{Cov}(X,Y)=E(XY)-EX\cdot EY=\frac83-\frac43\cdot2=0
\quad\Longrightarrow\quad \rho_{XY}=0,
\]
故 $X$ 与 $Y$ \textbf{不相关}。

\key{正交？\ 否}
正交的判据是 $E[XY]=0$，而 $E[XY]=\dfrac83\neq0$，故\textbf{不正交}。

\key{小结}
本题给出的是"独立、不相关、但\textbf{不正交}"的典型例。
原因：不相关 $\Leftrightarrow$ 正交 $\wedge$（$EX$ 或 $EY$ 为零）。本题 $EX=\tfrac43\neq0$ 且 $EY=2\neq0$，
所以"不相关"与"正交"并不等价。记忆要点：
\[
\text{独立}\;\Longrightarrow\;\text{不相关};\qquad
\text{不相关}+(\text{至少一个均值为 }0)\;\Longleftrightarrow\;\text{正交}.
\]
\end{tcolorbox}
\qtitle{第 6 题}{Day 2}
{\small\color{accentred}(Final examination 2024--25, Q1(b))}\par\smallskip
{\small\color{muted}教材小节：1.4.3 / 1.2.3}\par\medskip
\noindent\textbf{正交且不相关但不独立：最易考的三个概念辨析}\par\smallskip\begin{tcolorbox}[breakable,colback=qbg,colframe=qbg,boxrule=0pt]\[
X[n]=\sin(2\pi n+\phi),\qquad Y[n]=\cos(2\pi n+\phi),\qquad \phi\sim U(0,\pi).
\]

\key{第 1 步：化简}
$n$ 为整数时 $\sin2\pi n=0$、$\cos2\pi n=1$，由和角公式
\[
X[n]=\sin(2\pi n+\phi)=\sin\phi,\qquad Y[n]=\cos(2\pi n+\phi)=\cos\phi ,
\]
即两个信号都退化成 $\phi$ 的函数。$\phi$ 在 $(0,\pi)$ 上均匀分布，$f(\phi)=\dfrac1\pi$。

\key{第 2 步：一阶矩}
\[
E[X]=\frac1\pi\int_0^\pi\sin\phi\dd\phi=\frac1\pi\big[-\cos\phi\big]_0^\pi=\frac{2}{\pi}\approx0.6366,
\]
\[
E[Y]=\frac1\pi\int_0^\pi\cos\phi\dd\phi=\frac1\pi\big[\sin\phi\big]_0^\pi=0 .
\]

\key{第 3 步：互相关}
\[
E[XY]=\frac1\pi\int_0^\pi\sin\phi\cos\phi\dd\phi
=\frac{1}{2\pi}\int_0^\pi\sin2\phi\dd\phi
=\frac{1}{2\pi}\Big[-\frac{\cos2\phi}{2}\Big]_0^\pi=0 .
\]

\key{第 4 步：三个判定}
\begin{bul}
  \item \textbf{正交：成立。} $E[XY]=0$，满足正交定义。
  \item \textbf{不相关：成立。}
        \[
        \mathrm{Cov}(X,Y)=E[XY]-EX\cdot EY=0-\frac{2}{\pi}\cdot0=0 .
        \]
        这也与"不相关 $\Leftrightarrow$ 正交 $+$ 至少一个均值为零"一致：此处 $EY=0$。
  \item \textbf{独立：不成立。} 恒有
        \[
        X^2+Y^2=\sin^2\phi+\cos^2\phi\equiv1 ,
        \]
        两者之间存在确定性函数关系 $Y=\pm\sqrt{1-X^2}$。从密度看，$(X,Y)$ 的联合分布集中在
        单位圆这条二维零测集上，不可能等于 $f_Xf_Y$，故\textbf{不独立}。
\end{bul}

\key{小结}
本题是"正交、不相关，但\textbf{不独立}"的标准反例。
\textbf{只有在联合高斯分布下}，"不相关 $\Rightarrow$ 独立"才成立——
这正是命题人想考的边界条件。
\end{tcolorbox}
\qtitle{第 7 题}{Day 3}
{\small\color{accentred}(Textbook 1.14 / Required)}\par\smallskip
{\small\color{muted}教材小节：1.2.1 / 1.4.2}\par\medskip
\noindent\textbf{联合（累积）分布、边缘分布律与 $Z=XY$ 的分布}\par\smallskip\begin{tcolorbox}[breakable,colback=qbg,colframe=qbg,boxrule=0pt]联合分布律为 $p_{ij}=P(X=x_i,Y=y_j)$，全部 6 项之和 $=0.40+0.60=1$ \checkmark

\key{(1) 联合分布函数与联合密度}
用单位阶跃函数 $u(t)$（$t\ge0$ 时为 1，否则为 0）表示：
\[
F(x,y)=P(X\le x,\;Y\le y)=\sum_{x_i\le x}\sum_{y_j\le y}p_{ij}
\]
\[
\begin{aligned}
F(x,y)={}&0.07\,u(x)u(y+1)+0.18\,u(x)u(y)+0.15\,u(x)u(y-1)\\
&+0.08\,u(x-1)u(y+1)+0.32\,u(x-1)u(y)+0.20\,u(x-1)u(y-1).
\end{aligned}
\]
离散型随机变量的"密度函数"用冲激函数表示：
\[
\begin{aligned}
f(x,y)={}&0.07\delta(x)\delta(y+1)+0.18\delta(x)\delta(y)+0.15\delta(x)\delta(y-1)\\
&+0.08\delta(x-1)\delta(y+1)+0.32\delta(x-1)\delta(y)+0.20\delta(x-1)\delta(y-1).
\end{aligned}
\]

\key{(2) 边缘分布律}
对行求和得 $X$ 的分布律，对列求和得 $Y$ 的分布律：
\begin{center}
\renewcommand{\arraystretch}{1.3}
\small
\begin{tabular}{c|ccc|c}
\toprule
$X\backslash Y$ & $-1$ & $0$ & $1$ & $p_{X}$ \\
\midrule
$0$ & $0.07$ & $0.18$ & $0.15$ & $\mathbf{0.40}$ \\
$1$ & $0.08$ & $0.32$ & $0.20$ & $\mathbf{0.60}$ \\
\midrule
$p_{Y}$ & $\mathbf{0.15}$ & $\mathbf{0.50}$ & $\mathbf{0.35}$ & $1.00$ \\
\bottomrule
\end{tabular}
\end{center}

\key{(3) $Z=XY$ 的分布律}
$X$ 只取 $0$ 或 $1$：$X=0$ 时恒有 $Z=0$；$X=1$ 时 $Z=Y$。于是
\[
P(Z=0)=0.40+0.32=0.72,\qquad
P(Z=-1)=0.08,\qquad
P(Z=1)=0.20 .
\]
校验 $0.72+0.08+0.20=1$ \checkmark

\key{(4) 相关系数}
\[
EX=0(0.40)+1(0.60)=0.60,\qquad E(X^2)=0.60\ \Rightarrow\ D(X)=0.60-0.36=0.24,
\]
\[
EY=(-1)(0.15)+0(0.50)+1(0.35)=0.20,\qquad E(Y^2)=0.15+0.35=0.50\ \Rightarrow\ D(Y)=0.50-0.04=0.46,
\]
\[
E(XY)=(1)(-1)(0.08)+(1)(1)(0.20)=0.12,
\]
\[
\mathrm{Cov}(X,Y)=0.12-(0.60)(0.20)=0
\quad\Longrightarrow\quad \boxed{\rho_{XY}=0}
\]

\key{附注（重要）}
$\rho_{XY}=0$ 说明 $X,Y$ 不相关，但二者\textbf{并不独立}：例如
\[
P(X=0,Y=0)=0.18\neq0.40\times0.50=P(X=0)P(Y=0).
\]
这是又一个"不相关但不独立"的例子，与第 6 题互为印证。
\end{tcolorbox}
\qtitle{第 8 题}{Day 3}
{\small\color{accentred}(Final examination 2025--26, Q1(a)--(c))}\par\smallskip
{\small\color{muted}教材小节：1.2.1 / 1.4.2}\par\medskip
\noindent\textbf{由随机试验定义的信号：一维/二维密度与协方差}\par\smallskip\begin{tcolorbox}[breakable,colback=qbg,colframe=qbg,boxrule=0pt]\[
X(t)=\begin{cases}\sin(2\pi t), & \text{正面}\\[2pt] t^2, & \text{反面}\end{cases}
\qquad P(\text{正})=P(\text{反})=\tfrac12 .
\]
注意 $t=1$ 时 $\sin2\pi=0$、$1^2=1$；$t=2$ 时 $\sin4\pi=0$、$2^2=4$。

\key{(a) 一维密度 $f_X(x,2)$}
$t=2$ 时 $X(2)$ 以 $\tfrac12$ 取 $0$、以 $\tfrac12$ 取 $4$，故
\[
\boxed{f_X(x,2)=\tfrac12\delta(x)+\tfrac12\delta(x-4)}
\]

\key{(b) 二维密度 $f_X(x_1,x_2;1,2)$}
$t_1=1$：取 $0$（正）或 $1$（反）；$t_2=2$：取 $0$（正）或 $4$（反）。
\textbf{同一次抛掷同时决定两个时刻的取值}，所以 4 种组合里只有 2 种会出现：
\[
(\text{正})\;(0,0)\ \text{概率}\tfrac12;\qquad
(\text{反})\;(1,4)\ \text{概率}\tfrac12 .
\]
\[
\boxed{f_X(x_1,x_2;1,2)=\tfrac12\delta(x_1)\delta(x_2)+\tfrac12\delta(x_1-1)\delta(x_2-4)}
\]
（两个时刻完全相关，这正是它与第 9 题的根本区别。）

\key{(c) 协方差 $C_X(t_1,t_2)$}
引入等概率取 $\pm1$ 的随机变量 $C$，则 $X(t)$ 可写成"均值 $+$ 随机部分"：
\[
X(t)=m(t)+C\,d(t),\qquad
m(t)=\frac{\sin2\pi t+t^2}{2},\qquad
d(t)=\frac{\sin2\pi t-t^2}{2}.
\]
由于 $E[C]=0$、$E[C^2]=1$，
\[
E[X(t_1)X(t_2)]=m(t_1)m(t_2)+d(t_1)d(t_2),
\]
\[
\boxed{C_X(t_1,t_2)=\mathrm{Cov}\big(X(t_1),X(t_2)\big)=d(t_1)d(t_2)
=\frac14\big(\sin2\pi t_1-t_1^2\big)\big(\sin2\pi t_2-t_2^2\big)}
\]

\key{校验}
取 $t_1=t_2=t$：$C_X(t,t)=\tfrac14(\sin2\pi t-t^2)^2=D(X(t))$。
$X(t)$ 是等概率两点分布，方差确为 $\big(\tfrac{a-b}{2}\big)^2$ \checkmark\\[2pt]
取 $t_1=1,t_2=2$：$C_X(1,2)=\tfrac14(0-1)(0-4)=1$。
直接验证：$E[X(1)X(2)]=\tfrac12(0)+\tfrac12(1\cdot4)=2$，$E[X(1)]=0.5$，$E[X(2)]=2$，
$\mathrm{Cov}=2-0.5\times2=1$ \checkmark
\end{tcolorbox}
\qtitle{第 9 题}{Day 3}
{\small\color{accentred}(Final examination 2025--26, Q1(d))}\par\smallskip
{\small\color{muted}教材小节：1.2.1 / 1.2.3}\par\medskip
\noindent\textbf{独立重复试验：二维密度可分解}\par\smallskip\begin{tcolorbox}[breakable,colback=qbg,colframe=qbg,boxrule=0pt]\[
Y[n]=\begin{cases}\sin(2\pi n), & \text{第 } n \text{ 次抛掷为正面}\\[2pt] n^2, & \text{第 } n \text{ 次抛掷为反面}\end{cases}
\qquad\text{各次抛掷独立，}P=\tfrac12 .
\]
因为 $n$ 为整数，$\sin(2\pi n)=0$，所以
\[
P(Y[n]=0)=\tfrac12,\qquad P(Y[n]=n^2)=\tfrac12 .
\]

\key{一维密度}
\[
n=1:\quad f_{Y_1}(y_1)=\tfrac12\delta(y_1)+\tfrac12\delta(y_1-1)
\]
\[
n=2:\quad f_{Y_2}(y_2)=\tfrac12\delta(y_2)+\tfrac12\delta(y_2-4)
\]

\key{二维密度（独立性 $\Rightarrow$ 边缘之积）}
各次抛掷相互独立，故 $Y[1]$ 与 $Y[2]$ 独立，二维密度等于两个一维密度之积，展开得 4 个冲激项：
\[
\boxed{f_Y(y_1,y_2;1,2)=\tfrac14\Big[\delta(y_1)\delta(y_2)+\delta(y_1)\delta(y_2-4)
+\delta(y_1-1)\delta(y_2)+\delta(y_1-1)\delta(y_2-4)\Big]}
\]

\key{与第 8 题(b) 的对比（考点精髓）}
\begin{center}
\renewcommand{\arraystretch}{1.3}
\small
\begin{tabularx}{\textwidth}{@{}L{3.4cm} c X@{}}
\toprule
& \textbf{项数} & \textbf{含义} \\
\midrule
第 8 题 (b) 同一次抛掷 & 2 项 & 两次取值完全相关，二维密度不可分解 \\
第 9 题 独立重复抛掷 & 4 项 & 独立，二维密度 $=$ 边缘之积 \\
\bottomrule
\end{tabularx}
\end{center}
判断二维密度能否写成边缘之积，是区分"独立"与"不独立"最快的办法。
\end{tcolorbox}
\qtitle{第 10 题}{Day 4}
{\small\color{accentred}(Formula recall \textemdash\ \S1.5.1)}\par\smallskip
{\small\color{muted}教材小节：1.5.1}\par\medskip
\noindent\textbf{特征函数的定义、变换对、唯一性与三条性质}\par\smallskip\begin{tcolorbox}[breakable,colback=qbg,colframe=qbg,boxrule=0pt]\key{(a) 定义}
特征函数是 $e^{jvX}$ 的数学期望：
\[
\phi_X(v)=E\big[e^{jvX}\big]=
\begin{cases}
\displaystyle\int_{-\infty}^{+\infty}f_X(x)e^{jvx}\dd x, & \text{连续型}\\[8pt]
\displaystyle\sum_k P_k\,e^{jvx_k}, & \text{离散型}
\end{cases}
\]
其中 $j=\sqrt{-1}$，$v$ 为实变量。

\key{(b) 傅里叶变换对与唯一性定理}
\[
\phi_X(v)=\int_{-\infty}^{+\infty}f_X(x)e^{jvx}\dd x,
\qquad
f_X(x)=\frac{1}{2\pi}\int_{-\infty}^{+\infty}\phi_X(v)e^{-jvx}\dd v .
\]
\textbf{唯一性定理}：$f_X(x)$ 与 $\phi_X(v)$ 一一对应、互相唯一确定，
因此特征函数是描述随机变量概率特性的另一种完整方式。
（注意：$\phi_X(v)$ 与 $v$ 本身都没有直接的物理意义。）

\key{(c) 性质 1（独立和 $\Rightarrow$ 乘积）}
若 $X_1,X_2,\cdots,X_n$ 相互独立，则
\[
\phi_{X_1+X_2+\cdots+X_n}(v)=\prod_{i=1}^{n}\phi_{X_i}(v).
\]
证明要点：$E\big[e^{jv\sum X_i}\big]=E\big[\prod e^{jvX_i}\big]=\prod E\big[e^{jvX_i}\big]$。

\key{(d) 性质 2（线性变换）}
\[
\phi_{aX+b}(v)=e^{jvb}\,\phi_X(av).
\]

\key{(e) 性质 3（由特征函数求矩）}
若 $E|X|^r<\infty$（$r$ 为非负整数），则对一切 $k\le r$，$\phi_X(v)$ 的 $k$ 阶导数存在且连续，并且
\[
\boxed{\,E\big[X^k\big]=j^{-k}\,\phi_X^{(k)}(0)\,}
\]
特别地 $E[X]=-j\,\phi_X'(0)$，$E[X^2]=-\phi_X''(0)$。
\end{tcolorbox}
\qtitle{第 11 题}{Day 4}
{\small\color{accentred}(Final examination 2022--23, Q1(b))}\par\smallskip
{\small\color{muted}教材小节：1.5.1}\par\medskip
\noindent\textbf{均匀分布的特征函数与线性变换性质}\par\smallskip\begin{tcolorbox}[breakable,colback=qbg,colframe=qbg,boxrule=0pt]\[
X(t)=U+t,\qquad t\in[-2,2],\qquad U\sim U[0,\pi].
\]

\key{第 1 步：求 $U$ 的特征函数}
$U$ 的密度为 $f_U(u)=\dfrac1\pi$（$0\le u\le\pi$），于是
\[
\phi_U(v)=\int_0^\pi\frac1\pi e^{jvu}\dd u
=\frac{1}{jv\pi}\Big[e^{jvu}\Big]_0^\pi
=\frac{e^{jv\pi}-1}{jv\pi}.
\]

\key{第 2 步：用性质 2 平移}
$X(t)=U+t$，即 $a=1$、$b=t$。由第 10 题性质 2 的 $\phi_{aX+b}(v)=e^{jvb}\phi_X(av)$：
\[
\boxed{\phi_{X(t)}(v)=e^{jvt}\phi_U(v)=e^{jvt}\,\frac{e^{jv\pi}-1}{jv\pi}}
\]

\key{等价形式（更直观）}
把分子配方：
\[
\frac{e^{jv\pi}-1}{jv\pi}
=\frac{e^{jv\pi/2}\big(e^{jv\pi/2}-e^{-jv\pi/2}\big)}{jv\pi}
=e^{jv\pi/2}\cdot\frac{\sin(v\pi/2)}{v\pi/2},
\]
故
\[
\phi_{X(t)}(v)=e^{jv\left(t+\frac{\pi}{2}\right)}\cdot\frac{\sin(v\pi/2)}{v\pi/2}.
\]
相位项表明分布中心在 $t+\dfrac\pi2$（即 $EU+t$），幅度项是均匀分布特征的 $\mathrm{sinc}$ 形状。

\key{校验}
令 $v=0$ 应得 $\phi=1$。对原式用洛必达或直接展开：
$\dfrac{e^{jv\pi}-1}{jv\pi}\to\dfrac{jv\pi}{jv\pi}=1$ \checkmark
\end{tcolorbox}
\qtitle{第 12 题}{Day 4}
{\small\color{accentred}(Textbook 1.29 / Optional)}\par\smallskip
{\small\color{muted}教材小节：1.5.1}\par\medskip
\noindent\textbf{由特征函数反求密度函数与各阶矩}\par\smallskip\begin{tcolorbox}[breakable,colback=qbg,colframe=qbg,boxrule=0pt]\key{(1) $\phi(v)=0.2+0.3e^{j2v}+0.2e^{j4v}+0.2e^{-j2v}+0.1e^{-j4v}$}
系数之和 $=0.2+0.3+0.2+0.2+0.1=1$，故为离散型，指数上的系数即为取值：
\[
f(x)=0.2\delta(x)+0.3\delta(x-2)+0.2\delta(x-4)+0.2\delta(x+2)+0.1\delta(x+4).
\]
\[
E X=0(0.2)+2(0.3)+4(0.2)+(-2)(0.2)+(-4)(0.1)=0.6,
\]
\[
E X^2=0+4(0.3)+16(0.2)+4(0.2)+16(0.1)=6.8,
\qquad
D(X)=6.8-0.36=\boxed{6.44}
\]

\key{(2) $\phi(v)=0.3e^{jv}+0.7e^{-jv}$}
$X$ 取 $1$ 的概率 $0.3$，取 $-1$ 的概率 $0.7$：
\[
f(x)=0.3\delta(x-1)+0.7\delta(x+1),
\]
\[
EX=0.3-0.7=-0.4,\qquad EX^2=1,\qquad D(X)=1-0.16=\boxed{0.84}
\]

\key{(3) $\phi(v)=\dfrac{4}{4-jv}=\dfrac{1}{1-jv/4}$}
与指数分布的特征函数 $\dfrac{\lambda}{\lambda-jv}=\dfrac{1}{1-jv/\lambda}$ 对照，得 $\lambda=4$：
\[
X\sim\mathrm{Exp}(4),\qquad f(x)=4e^{-4x}\ (x\ge0),
\]
\[
EX=\frac1\lambda=\frac14=0.25,\qquad
EX^2=\frac{2}{\lambda^2}=\frac{1}{8}=0.125,\qquad
D(X)=\frac{1}{\lambda^2}=\boxed{0.0625}
\]

\key{(4) $\phi(v)=\dfrac{\sin 5v}{5v}$}
由 $U(a,b)$ 的特征函数 $\dfrac{e^{jvb}-e^{jva}}{jv(b-a)}$，代入 $a=-5,b=5$：
\[
\frac{e^{j5v}-e^{-j5v}}{jv\cdot10}=\frac{2j\sin5v}{10jv}=\frac{\sin5v}{5v},
\]
故 $X\sim U(-5,5)$：
\[
f(x)=\frac{1}{10}\ (|x|<5),\qquad
EX=0,\qquad
EX^2=\frac{(b-a)^2}{12}=\frac{100}{12}=\boxed{\frac{25}{3}\approx8.333},\qquad
D(X)=\frac{25}{3}
\]
\end{tcolorbox}
\qtitle{第 13 题}{Day 5}
{\small\color{accentred}(Formula recall \textperiodcentered\ \S1.6)}\par\smallskip
{\small\color{muted}教材小节：1.6}\par\medskip
\noindent\textbf{六个典型分布一览（必背）}\par\smallskip\begin{tcolorbox}[breakable,colback=qbg,colframe=qbg,boxrule=0pt]\begin{center}
\renewcommand{\arraystretch}{1.75}
\scriptsize
\begin{tabularx}{\textwidth}{@{}c L{2.5cm} X X X@{}}
\toprule
& \textbf{分布} & \textbf{密度 / 分布律} & \textbf{特征函数 $\phi_X(v)$} & \textbf{均值、方差} \\
\midrule
(a) & 均匀 $U(a,b)$ & $f=\dfrac{1}{b-a},\ a<x<b$
    & $\dfrac{e^{jvb}-e^{jva}}{jv(b-a)}$
    & $EX=\dfrac{a+b}{2}$；$D(X)=\dfrac{(b-a)^2}{12}$ \\
(b) & 指数（参数 $\lambda$） & $f=\lambda e^{-\lambda x},\ x\ge0$
    & $\dfrac{\lambda}{\lambda-jv}$
    & $EX=\dfrac1\lambda$；$D(X)=\dfrac{1}{\lambda^2}$ \\
(c) & 高斯 $N(\mu,\sigma^2)$ & $f=\dfrac{1}{\sqrt{2\pi}\sigma}e^{-\frac{(x-\mu)^2}{2\sigma^2}}$
    & $e^{jv\mu-\frac12\sigma^2v^2}$
    & $EX=\mu$；$D(X)=\sigma^2$ \\
(d) & $0$--$1$ 分布（参数 $p$） & $P(X{=}1)=p,\ P(X{=}0)=q$
    & $q+pe^{jv}$
    & $EX=p$；$D(X)=pq$ \\
(e) & 二项 $B(n,p)$ & $P(X{=}k)=\dbinom{n}{k}p^kq^{n-k}$
    & $(q+pe^{jv})^n$
    & $EX=np$；$D(X)=npq$ \\
(f) & 泊松 $P(\lambda)$ & $P(X{=}k)=\dfrac{\lambda^k}{k!}e^{-\lambda}$
    & $e^{\lambda(e^{jv}-1)}$
    & $EX=\lambda$；$D(X)=\lambda$ \\
\bottomrule
\end{tabularx}
\end{center}

\vspace{2pt}
\key{记忆线索}
\begin{bul}
  \item 高斯、二项、泊松的特征函数都写成"指数型"：高斯指数里含 $v^2$，泊松指数里含 $e^{jv}$，
        二项是 $n$ 次幂——把 $0$--$1$ 分布的特征函数取 $n$ 次幂即得二项（由性质 1）。
  \item 均匀分布的密度是"矩形"，其特征函数是 $\mathrm{sinc}$ 型；
        指数分布的特征函数是分式 $\dfrac{\lambda}{\lambda-jv}$。
  \item 均值与方差可直接由定义记住：均匀取中点、指数是 $\lambda^{-1}$ 与 $\lambda^{-2}$、
        高斯是 $\mu$ 与 $\sigma^2$、二项是 $np$ 与 $npq$、泊松两者都是 $\lambda$。
\end{bul}
\end{tcolorbox}
\qtitle{第 14 题}{Day 5}
{\small\color{accentred}(Textbook Exercise 1.17)}\par\smallskip
{\small\color{muted}教材小节：1.3.1}\par\medskip
\noindent\textbf{非单调函数的密度：连续部分 $+$ 原点冲激}\par\smallskip\begin{tcolorbox}[breakable,colback=qbg,colframe=qbg,boxrule=0pt]\[
Y=g(X)=\begin{cases}X, & X\ge0\\ 0, & X<0\end{cases}
\]

\key{方法：先求分布函数 $F_Y(y)$，再求导}
\noindent\textbf{(i) $y<0$：} $Y\ge0$ 恒成立，故
\[
F_Y(y)=P(Y\le y)=0 .
\]
\noindent\textbf{(ii) $y\ge0$：} 关键在于化简事件。若 $X<0$ 则 $Y=0\le y$；
若 $0\le X\le y$ 则 $Y=X\le y$；反之若 $Y\le y$ 必有 $X\le y$。因此
\[
\{Y\le y\}=\{X\le y\}\quad(y\ge0)
\qquad\Longrightarrow\qquad
F_Y(y)=F_X(y)\ (y\ge0).
\]
合起来即
\[
F_Y(y)=F_X(y)\,u(y),
\]
其中 $u(y)$ 为单位阶跃函数。

\key{求导得密度}
对乘积求导，并注意 $F_X(y)\delta(y)=F_X(0)\delta(y)$：
\[
f_Y(y)=\frac{\dd F_Y}{\dd y}=f_X(y)u(y)+F_X(y)\delta(y),
\]
\[
\boxed{\,f_Y(y)=f_X(y)\,u(y)+F_X(0)\,\delta(y)\,}
\]

\key{物理含义（考试常问）}
\begin{bul}
  \item 连续部分：把 $f_X$ 的\textbf{正半轴}原样保留（$X\ge0$ 时 $Y=X$）。
  \item 冲激部分：整个负半轴 $\{X<0\}$ 被"压"到 $y=0$ 一点，
        其强度等于该事件的概率 $P(X\le0)=F_X(0)$。
\end{bul}
这就是"非单调变换会产生离散分量"的标准范例，也是 §1.3.1 的核心技巧。
\end{tcolorbox}
\qtitle{第 15 题}{Day 5}
{\small\color{accentred}(Adapted from 2023--24 Q2(c))}\par\smallskip
{\small\color{muted}教材小节：1.6 / 1.4.2}\par\medskip
\noindent\textbf{白噪声与高斯分布的条件辨析}\par\smallskip\begin{tcolorbox}[breakable,colback=qbg,colframe=qbg,boxrule=0pt]\key{(1) 只给白噪声，不能推出高斯}
白噪声规定二阶谱为常数，没有规定全部有限维分布。线性系统保持已有高斯性，却不会把任何非高斯输入自动变成精确高斯。因此原条件只能确定$Y$的均值和方差，不能唯一给出其PDF。
例如取高斯白噪声$G$及独立随机幅度$A$，令$P(A=0)=P(A=\sqrt2)=1/2$，$W=AG$。因$E(A^2)=1$，$W$仍具有相同白噪声谱；滤波差分输出却一半概率恒为0，是高斯尺度混合而非非退化高斯。这是按二阶白噪声定义的反例。
\key{(2) 另加高斯假设，才能使用§2.4.3}
滤波和延时差分都是线性操作，故$Y(t)$高斯，且$m_Y=0$。由相关展开，
\[\begin{aligned}v_Y&=E[(X(t)-X(t-T))^2]\\&=2R_X(0)-2R_X(T)\\&=\boxed{\frac{N_0}{2}(1-e^{-T})}>0.\end{aligned}\]
于是
\[\boxed{f_Y(y;t)=\frac1{\sqrt{\pi N_0(1-e^{-T})}}\exp\!\left[-\frac{y^2}{N_0(1-e^{-T})}\right].}\]
注意密度指数的分母是$2v_Y$，归一化系数为$1/\sqrt{2\pi v_Y}$。
\key{(3) $T=0$}
此时$Y(t)\equiv0$，分布是$\delta(y)$，无普通一维密度，不能把$T=0$硬代入上面的非退化公式。

\end{tcolorbox}
\qtitle{第 16 题}{Day 6}
{\small\color{accentred}(Textbook 1.16 / Required)}\par\smallskip
{\small\color{muted}教材小节：1.3.2}\par\medskip
\noindent\textbf{线性变换求二维密度并用两种方法判独立}\par\smallskip\begin{tcolorbox}[breakable,colback=qbg,colframe=qbg,boxrule=0pt]$X\sim N(0,1)$、$Y\sim N(0,1)$ 且相互独立，$U=X+Y$，$V=X-Y$。

\key{(1) 联合密度——方法一：雅可比变换}
反解：$X=\dfrac{U+V}{2}$，$Y=\dfrac{U-V}{2}$。雅可比行列式
\[
J=\frac{\partial(x,y)}{\partial(u,v)}
=\begin{vmatrix} \tfrac12 & \tfrac12\\[2pt] \tfrac12 & -\tfrac12\end{vmatrix}
=-\frac14-\frac14=-\frac12,
\qquad |J|=\frac12 .
\]
因为 $X,Y$ 独立，
\[
f_{UV}(u,v)=f_X\!\Big(\frac{u+v}{2}\Big)f_Y\!\Big(\frac{u-v}{2}\Big)|J|
=\frac{1}{\sqrt{2\pi}}e^{-\frac{(u+v)^2}{8}}\cdot\frac{1}{\sqrt{2\pi}}e^{-\frac{(u-v)^2}{8}}\cdot\frac12 .
\]
把两个指数合并：$\dfrac{(u+v)^2+(u-v)^2}{8}=\dfrac{2u^2+2v^2}{8}=\dfrac{u^2+v^2}{4}$，于是
\[
\boxed{\,f_{UV}(u,v)=\frac{1}{4\pi}\exp\!\Big(-\frac{u^2+v^2}{4}\Big)\,}
\]

\key{(2) 是否独立——看能否分解}
把上式拆成两个因子：
\[
f_{UV}(u,v)=\underbrace{\frac{1}{\sqrt{4\pi}}e^{-u^2/4}}_{f_U(u)}\cdot
\underbrace{\frac{1}{\sqrt{4\pi}}e^{-v^2/4}}_{f_V(v)},
\qquad\text{且}\ \Big(\tfrac{1}{\sqrt{4\pi}}\Big)^2=\frac{1}{4\pi}\ \text{\checkmark}
\]
即 $U\sim N(0,2)$、$V\sim N(0,2)$ 且\textbf{相互独立}。

\key{方法二：联合高斯 $+$ 不相关 $\Rightarrow$ 独立（更省力）}
$(U,V)$ 是独立高斯变量的线性变换，故服从\textbf{联合高斯}分布。只需算协方差：
\[
\mathrm{Cov}(U,V)=E[UV]-EU\cdot EV=E\big[(X+Y)(X-Y)\big]-0=E[X^2]-E[Y^2]=1-1=0 .
\]
联合高斯且协方差为零 $\Rightarrow$ 不相关 $\Rightarrow$ \textbf{独立}；
且 $D(U)=D(X)+D(Y)=2$，$D(V)=2$。
两种方法都指向 $U,V\overset{\text{i.i.d.}}{\sim}N(0,2)$。
\end{tcolorbox}
\qtitle{第 17 题}{Day 6}
{\small\color{accentred}(Formula recall \textperiodcentered\ \S1.2.1 and \S1.6)}\par\smallskip
{\small\color{muted}教材小节：1.2.1 / 1.6}\par\medskip
\noindent\textbf{二维高斯密度与"两时刻"联合密度}\par\smallskip\begin{tcolorbox}[breakable,colback=qbg,colframe=qbg,boxrule=0pt]\key{(a) 一般二维高斯联合密度}
\[
f_{XY}(x,y)=\frac{1}{2\pi\sigma_X\sigma_Y\sqrt{1-\rho^2}}
\exp\Big\{-\frac{1}{2(1-\rho^2)}\Big[
\frac{(x-\mu_X)^2}{\sigma_X^2}
-\frac{2\rho(x-\mu_X)(y-\mu_Y)}{\sigma_X\sigma_Y}
+\frac{(y-\mu_Y)^2}{\sigma_Y^2}\Big]\Big\}
\]

\key{(b) 独立的充要条件：$\rho=0$}
当 $\rho=0$ 时，$1-\rho^2=1$，方括号中的交叉项消失：
\[
f_{XY}(x,y)=\frac{1}{\sqrt{2\pi}\sigma_X}e^{-\frac{(x-\mu_X)^2}{2\sigma_X^2}}
\cdot\frac{1}{\sqrt{2\pi}\sigma_Y}e^{-\frac{(y-\mu_Y)^2}{2\sigma_Y^2}}
=f_X(x)f_Y(y).
\]
联合密度可分解为边缘密度之积，故独立。
反过来，对\textbf{联合高斯}变量，独立 $\Rightarrow$ 不相关 $\Rightarrow\rho=0$，
所以 $\rho=0$ 是独立与不相关的\textbf{共同}充要条件。
（提醒：这只对联合高斯成立——第 5、7 题就是反例。）

\key{(c) 特例 $\sigma_X=\sigma_Y=\sigma,\ \rho=0$}
\[
\boxed{\,f_{XY}(x,y)=\frac{1}{2\pi\sigma^2}
\exp\Big\{-\frac{(x-\mu_X)^2+(y-\mu_Y)^2}{2\sigma^2}\Big\}\,}
\]
指数只与到中心 $(\mu_X,\mu_Y)$ 的距离有关，是\textbf{圆对称（各向同性）}的高斯密度，
等高线为同心圆。窄带噪声的同相/正交分量正是这种情形。

\key{(d) 高斯过程的二维密度由什么决定}
高斯过程的一切有限维分布都是高斯的，而高斯分布由\textbf{一阶矩和二阶矩}完全确定。
因此 $f_X(x_1,x_2;t_1,t_2)$ 只由两个量决定：
\begin{enumerate}[label=(\roman*),leftmargin=1.7em,itemsep=1pt]
  \item \textbf{均值} $m_X(t)$，
  \item \textbf{自协方差} $C_X(t_1,t_2)$（或等价地用自相关 $R_X$ 与均值）。
\end{enumerate}
写成矩阵形式（$\mathbf{x}=(x_1,x_2)^{\mathrm T}$，$\mathbf{m}=\big(m_X(t_1),m_X(t_2)\big)^{\mathrm T}$）：
\[
f_X(x_1,x_2;t_1,t_2)=\frac{1}{2\pi\sqrt{\det K}}
\exp\Big\{-\frac12(\mathbf{x}-\mathbf{m})^{\mathrm T}K^{-1}(\mathbf{x}-\mathbf{m})\Big\},
\]
\[
K=\begin{pmatrix}\sigma_1^2 & C_X(t_1,t_2)\\ C_X(t_1,t_2) & \sigma_2^2\end{pmatrix},
\qquad \sigma_i^2=C_X(t_i,t_i).
\]
\textbf{平稳情形}：$m_X$ 为常数，$C_X(t_1,t_2)=C_X(\tau)$（$\tau=t_2-t_1$），
$\sigma_1^2=\sigma_2^2=C_X(0)$，相关系数 $\rho_\tau=\dfrac{C_X(\tau)}{C_X(0)}$。
第 18 题(c) 正是这个公式的直接应用。
\key{普通密度存在的条件}
上面的普通二维密度要求两方差为正且$|\rho|<1$，等价于$K$正定。若$K$奇异，高斯分布仍可由均值和协方差确定，但集中在低维集合，没有普通二维PDF；此时不能使用$K^{-1}$或$1/\sqrt{\det K}$。

\end{tcolorbox}
\qtitle{第 18 题}{Day 6}
{\small\color{accentred}(Final examination 2025--26, Q4(b)--(c))}\par\smallskip
{\small\color{muted}教材小节：1.6 / 1.4.2}\par\medskip
\noindent\textbf{窄带高斯同相/正交分量的联合密度}\par\smallskip\begin{tcolorbox}[breakable,colback=qbg,colframe=qbg,boxrule=0pt]$X(t)=i(t)\cos\omega_0t-q(t)\sin\omega_0t$，$\omega_0=148\pi$，
$S_X(\omega)=1$（$144\pi\le|\omega|\le156\pi$），其余为 $0$。

\key{预备：自相关与两个基本关系}
\[
R_X(\tau)=\frac{1}{2\pi}\int_{-\infty}^{+\infty}\!\!S_X(\omega)e^{j\omega\tau}\dd\omega
=\frac{1}{\pi}\int_{144\pi}^{156\pi}\!\!\cos(\omega\tau)\dd\omega
=\frac{\sin(156\pi\tau)-\sin(144\pi\tau)}{\pi\tau},
\]
故
\[
R_X(0)=\lim_{\tau\to0}\frac{(156\pi-144\pi)\tau}{\pi\tau}=12
\quad\Longrightarrow\quad
\boxed{\sigma_i^2=\sigma_q^2=R_i(0)=R_X(0)=12}
\]
另一关键关系（$R_i(\tau)=R_q(\tau)=R_X\cos\omega_0\tau+\hat R_X\sin\omega_0\tau$）：
\[
\boxed{\,R_{iq}(\tau)=\hat R_X(\tau)\cos\omega_0\tau-R_X(\tau)\sin\omega_0\tau\,}
\]
其中 $\hat R_X$ 为 $R_X$ 的希尔伯特变换，对本例
\[
\hat R_X(\tau)=\frac{\cos(144\pi\tau)-\cos(156\pi\tau)}{\pi\tau}.
\]

\key{(a) 同一时刻：$f_{iq}(i,q;t,t)$}
此时 $\tau=0$：$\hat R_X(0)=0$（奇函数），故
\[
R_{iq}(0)=0\times1-12\times0=0 .
\]
$i,q$ 零均值、联合高斯、互不相关 $\Rightarrow$ \textbf{相互独立}，且各自方差为 $12$：
\[
\boxed{\,f_{iq}(i,q;t,t)=\frac{1}{2\pi\cdot12}\exp\!\Big(-\frac{i^2+q^2}{2\cdot12}\Big)
=\frac{1}{24\pi}\exp\!\Big(-\frac{i^2+q^2}{24}\Big)\,}
\]

\key{(b) 时延 $\tau=\tfrac{1}{12}$：三处计算（命题人精心设计的数值）}
\[
\omega_0\tau=\frac{148\pi}{12}=12\pi+\frac{\pi}{3}
\ \Longrightarrow\ \cos\omega_0\tau=\frac12,\qquad \sin\omega_0\tau=\frac{\sqrt3}{2};
\]
\[
R_X\!\Big(\tfrac1{12}\Big)=\frac{\sin 13\pi-\sin 12\pi}{\pi/12}=0 ;
\]
\[
\hat R_X\!\Big(\tfrac1{12}\Big)=\frac{\cos 12\pi-\cos 13\pi}{\pi/12}
=\frac{1-(-1)}{\pi/12}=\frac{24}{\pi}.
\]
于是
\[
R_{iq}\!\Big(\tfrac1{12}\Big)=\frac{24}{\pi}\cdot\frac12-0\cdot\frac{\sqrt3}{2}=\frac{12}{\pi},
\qquad
\rho=\frac{R_{iq}(1/12)}{\sigma_i\sigma_q}=\frac{12/\pi}{12}=\frac1\pi .
\]

\key{(b) 结果：相关系数为 $1/\pi$ 的二维高斯密度}
两个边缘均为 $N(0,12)$，相关系数 $\rho=1/\pi$，代入第 17 题(a) 的公式：
\[
\boxed{\,f_{iq}\!\Big(i,q;t,t+\tfrac1{12}\Big)
=\frac{1}{24\pi\sqrt{1-\pi^{-2}}}\exp\!\Big\{-\frac{i^2-\frac{2}{\pi}iq+q^2}{24\big(1-\pi^{-2}\big)}\Big\}\,}
\]
数值形式：
\[
f_{iq}\approx0.01399\,\exp\!\Big\{-\frac{i^2-0.6366\,iq+q^2}{21.57}\Big\}
\]

\key{为什么要小心（本题的"坑"）}
若沿用"窄带噪声的 $i$、$q$ 恒不相关"这一\textbf{只在频谱关于 $\omega_0$ 对称时成立}的结论，
就会得到与 (a) 完全相同的答案。但本题频谱为 $[144\pi,156\pi]$，中心是 $150\pi\neq\omega_0=148\pi$，
\textbf{并不关于 $\omega_0$ 对称}，因此 $R_{iq}(\tau)\not\equiv0$，交叉项必须保留。
$\tau=\tfrac1{12}$ 恰好让 $R_X(1/12)=0$、$\cos\omega_0\tau=\tfrac12$，
使交叉项只由希尔伯特变换项贡献——数值如此"整齐"，正是命题人有意为之。
\end{tcolorbox}
\qtitle{第 19 题}{Day 7}
{\small\color{accentred}(Textbook 1.6 / Required)}\par\smallskip
{\small\color{muted}教材小节：1.1.1}\par\medskip
\noindent\textbf{全概率公式与贝叶斯公式}\par\smallskip\begin{tcolorbox}[breakable,colback=qbg,colframe=qbg,boxrule=0pt]设 $B_i$ 表示"取到第 $i$ 批"，$i=1,2,3,4$；$D$ 表示"取到次品"。
\textbf{关键是"随机选一个批次"}，所以批次先验概率均等，与各批数量无关：
\[
P(B_1)=P(B_2)=P(B_3)=P(B_4)=\frac14 .
\]
条件次品率：
\[
P(D|B_1)=0.05,\quad P(D|B_2)=0.40,\quad P(D|B_3)=P(D|B_4)=0.10 .
\]

\key{(1) 全概率公式}
\[
P(D)=\sum_{i=1}^{4}P(B_i)P(D|B_i)
=\frac14\big(0.05+0.40+0.10+0.10\big)=\frac{0.65}{4}=0.1625 .
\]

\key{(2) 贝叶斯公式}
\[
P(B_2|D)=\frac{P(D|B_2)P(B_2)}{P(D)}
=\frac{0.40\times0.25}{0.1625}=\frac{0.10}{0.1625}=\frac{8}{13}\approx0.6154 .
\]

\key{易错点}
\begin{bul}
  \item 四批的批量 $2000$、$500$、$1000$、$1000$ 是\textbf{干扰信息}。
        若误按"从全部零件中随机抽一个"，则先验应取 $P(B_i)=n_i/\sum n_i$，答案完全不同。
  \item 结论的合理性检查：次品率最高的第二批（$40\%$）应当占据大部分后验概率，
        $0.6154$ 大于其先验 $0.25$，符合贝叶斯更新方向。
\end{bul}
\end{tcolorbox}
\qtitle{第 20 题}{Day 7}
{\small\color{accentred}(Textbook 1.24 / Required)}\par\smallskip
{\small\color{muted}教材小节：1.4.4}\par\medskip
\noindent\textbf{条件数学期望与重期望公式}\par\smallskip\begin{tcolorbox}[breakable,colback=qbg,colframe=qbg,boxrule=0pt]$X\sim U[0,a]$；给定 $X$ 后，$Y$ 在 $[X,a]$ 上均匀分布。

\key{(1) 条件期望 $E(Y|X)$}
给定 $X$ 时，$Y$ 是区间 $[X,a]$ 上的均匀分布，取区间中点：
\[
\boxed{\,E(Y|X)=\frac{X+a}{2},\qquad 0\le X\le a\,}
\]

\key{(2) 无条件期望 $EY$——用重期望（全期望）公式}
\[
E[Y]=E\big\{E[Y|X]\big\}=E\Big[\frac{X+a}{2}\Big]=\frac{E[X]+a}{2}.
\]
由 $X\sim U[0,a]$ 得 $E[X]=\dfrac a2$，于是
\[
\boxed{\,E[Y]=\frac{\frac a2+a}{2}=\frac{3a}{4}\,}
\]

\key{另解（顺带验证）}
先求 $Y$ 的边缘密度。由 $f_{Y|X}(y|x)=\dfrac{1}{a-x}$（$x<y<a$）与 $f_X(x)=\dfrac1a$：
\[
f_{XY}(x,y)=\frac{1}{a(a-x)},\qquad
f_Y(y)=\int_0^y\frac{\dd x}{a(a-x)}=\frac1a\ln\frac{a}{a-y}\quad(0<y<a).
\]
再算 $E[Y]=\displaystyle\int_0^a y\cdot\frac1a\ln\frac{a}{a-y}\dd y=\frac{3a}{4}$，与上面一致。
\textbf{但显然重期望公式省力得多}——条件期望类题目应优先考虑它。
\end{tcolorbox}
\qtitle{第 21 题}{Day 7}
{\small\color{accentred}(Textbook 1.33 / Required)}\par\smallskip
{\small\color{muted}教材小节：1.5.2}\par\medskip
\noindent\textbf{由联合特征函数求边缘特征函数与矩}\par\smallskip\begin{tcolorbox}[breakable,colback=qbg,colframe=qbg,boxrule=0pt]\[
\phi_{XY}(u,v)=\frac{6}{6-2ju-3jv-uv}.
\]

\key{(1) $X$ 的特征函数：令 $v=0$}
\[
\phi_X(u)=\phi_{XY}(u,0)=\frac{6}{6-2ju}=\frac{3}{3-ju}=\frac{1}{1-ju/3}.
\]
与 $\dfrac{\lambda}{\lambda-ju}=\dfrac{1}{1-ju/\lambda}$ 对照得 $\lambda=3$，故
\[
\boxed{X\sim\mathrm{Exp}(3),\qquad f_X(x)=3e^{-3x}\ (x\ge0)}
\]

\key{(2) $Y$ 的期望与方差：令 $u=0$}
\[
\phi_Y(v)=\phi_{XY}(0,v)=\frac{6}{6-3jv}=\frac{2}{2-jv}=\frac{1}{1-jv/2}
\quad\Longrightarrow\quad Y\sim\mathrm{Exp}(2).
\]
\[
\boxed{E[Y]=\frac{1}{\lambda}=\frac12,\qquad D(Y)=\frac{1}{\lambda^2}=\frac14}
\]

\key{用性质 3 复核（矩发生性质 $E[Y^k]=j^{-k}\phi_Y^{(k)}(0)$）}
令 $\phi_Y(v)=(1-jv/2)^{-1}$：
\[
\phi_Y'(v)=\frac{j}{2}\Big(1-\frac{jv}{2}\Big)^{-2}
\ \Rightarrow\ \phi_Y'(0)=\frac j2
\ \Rightarrow\ E[Y]=j^{-1}\cdot\frac j2=\frac12\ \text{\checkmark}
\]
\[
\phi_Y''(v)=-\frac12\Big(1-\frac{jv}{2}\Big)^{-3}
\ \Rightarrow\ \phi_Y''(0)=-\frac12
\ \Rightarrow\ E[Y^2]=j^{-2}\Big(-\frac12\Big)=\frac12\ \text{\checkmark}
\]
\[
D(Y)=E[Y^2]-(E[Y])^2=\frac12-\frac14=\frac14\ \text{\checkmark}
\]

\key{方法要点}
\textbf{求某个变量的边缘特征函数，就把联合特征函数中另一个变量置零}。
这是 §1.5.2 最常考的一个动作。
\end{tcolorbox}
\qtitle{第 22 题}{Day 8}
{\small\color{accentred}(2018--19 Q1 (201812))}\par\smallskip
{\small\color{muted}教材小节：1.3.1 / 1.6}\par\medskip
\noindent\textbf{均匀变量的仿射变换}\par\smallskip\begin{tcolorbox}[breakable,colback=qbg,colframe=qbg,boxrule=0pt]\key{先确定支持域（§1.3.1）}
$0\le X\le1$映射为$1\le Y\le6$，逆变换$x=(y-1)/5$，$|dx/dy|=1/5$。
\[\boxed{f_Y(y)=\begin{cases}1/5,&1\le y\le6,\\0,&\text{其他}.\end{cases}}\]
连续分布在端点上的密度取值不影响概率。
\key{用CDF独立复核}
$y<1$时$F_Y(y)=0$；$1\le y<6$时$F_Y(y)=P(X\le(y-1)/5)=(y-1)/5$；$y\ge6$时$F_Y(y)=1$。对中段求导仍为$1/5$。
\key{归一化与矩}
$\int_1^6dy/5=1$。作为附加检查，$E(Y)=5/2+1=7/2$，$\operatorname{Var}(Y)=25/12$。变换后的区间长度是5，密度不是原来的1。
\end{tcolorbox}
\qtitle{第 23 题}{Day 8}
{\small\color{accentred}(Adapted from 2019--20 Q1 (201912))}\par\smallskip
{\small\color{muted}教材小节：1.3.1 / 1.5.1 / 1.6}\par\medskip
\noindent\textbf{标准正态与截断正态必须区分}\par\smallskip\begin{tcolorbox}[breakable,colback=qbg,colframe=qbg,boxrule=0pt]\key{明确标准正态假设下的答案}
由§1.3.1，$Y\sim N(2,36)$，所以
\[\boxed{f_Y(y)=\frac{1}{6\sqrt{2\pi}}e^{-(y-2)^2/72},\quad y\in\mathbb R.}\]
由§1.5.1，$\phi_Y(v)=e^{j2v}\phi_X(6v)=\boxed{e^{j2v-18v^2}}$。在$v=0$为1，二阶累积量为36。
\key{原题为什么不能直接写唯一答案}
普通非退化正态变量没有有限区间支持。“in [0,1]”既可能想表达截断，也可能误把均值0、方差1写成区间；原文未澄清。上半题额外指定$N(0,1)$。
\key{若明确指定截断标准正态}
记标准正态密度$\varphi$与CDF$\Phi$，$Z=\Phi(1)-\Phi(0)>0$。条件密度为$\varphi(x)/Z$，$0\le x\le1$。
\[f_Y(y)=\begin{cases}\dfrac{\varphi((y-2)/6)}{6Z},&2\le y\le8,\\0,&\text{其他},\end{cases}\]
\[\phi_Y(v)=\frac{e^{j2v}}{Z}\int_0^1 e^{j6vx}\varphi(x)\,dx.\]
此时不能用$e^{j2v-18v^2}$；代入$v=0$可核对条件密度归一化。
\end{tcolorbox}
\qtitle{第 24 题}{Day 8}
{\small\color{accentred}(2019--20 Q2 (201912))}\par\smallskip
{\small\color{muted}教材小节：1.2.1}\par\medskip
\noindent\textbf{指数型联合密度的边缘化}\par\smallskip\begin{tcolorbox}[breakable,colback=qbg,colframe=qbg,boxrule=0pt]\key{边缘密度是积分，不是固定切片（§1.2.1）}
对$x>0$，
\[f_X(x)=\int_0^\infty xe^{-x(y+1)}dy=xe^{-x}\frac1x=\boxed{e^{-x}}.\]
对$y\ge0$，利用$\int_0^\infty xe^{-cx}dx=1/c^2$，得
\[\boxed{f_Y(y)=\frac1{(y+1)^2}}.\]
区间外均为0。密度在$x=0$的单点取值不影响分布，不应因原式$x=0$为0而误删指数边缘。
\key{归一化}
$\int_0^\infty e^{-x}dx=1$，$\int_0^\infty (1+y)^{-2}dy=1$，故联合密度亦归一化。
\key{独立性（§1.2.3）}
$f_X(x)f_Y(y)=e^{-x}/(1+y)^2$与原联合密度一般不相等，故不独立。
在$x>0$，$f_{Y|X}(y|x)=xe^{-xy}$，即给定$X=x$后$Y$为参数$x$的指数变量，清楚显示依赖。延伸：$E(Y)=\infty$，不能未经检查套用有限矩公式。
\end{tcolorbox}
\qtitle{第 25 题}{Day 9}
{\small\color{accentred}(Adapted from 2020--21 Q4(b,c) (202104))}\par\smallskip
{\small\color{muted}教材小节：1.6 / 2.4.3}\par\medskip
\noindent\textbf{滤波输出的高斯密度与对称概率}\par\smallskip\begin{tcolorbox}[breakable,colback=qbg,colframe=qbg,boxrule=0pt]\key{先从相关函数提取方差}
零均值输入经LTI仍零均值。于是
\[v=R_Y(0)=\frac{\sigma_X^2(1-a/b)}{(b-a)(b+a)}=\boxed{\frac{\sigma_X^2}{b(a+b)}}>0.\]
不要把指数相关中的$a$或滤波器中的$b$直接当作标准差。
\key{高斯性与密度}
线性滤波保持输入已给定的高斯性（§2.4.3）。仅知道二阶相关并不能证明高斯。本题明确给Gaussian，故由§1.6，
\[\boxed{f_Y(y;t)=\frac{1}{\sqrt{2\pi v}}e^{-y^2/(2v)},\quad y\in\mathbb R.}\]
均值为0且$v>0$，分布连续并关于原点对称，故$P(Y=0)=0$，
\[\boxed{P\{Y(t)\ge0\}=\tfrac12}.\]
若输入方差为0，则退化为恒0，概率为1；这就是非退化条件不可省略的原因。
\end{tcolorbox}
\qtitle{第 26 题}{Day 9}
{\small\color{accentred}(2021--22 Q1(c) (202112))}\par\smallskip
{\small\color{muted}教材小节：1.5.2 / 1.4.2}\par\medskip
\noindent\textbf{两时刻高斯向量的联合特征函数}\par\smallskip\begin{tcolorbox}[breakable,colback=qbg,colframe=qbg,boxrule=0pt]\key{固定时刻后先写出随机变量}
$Z_1=A+B$，$Z_2=A$，均值都为0。独立给$E(AB)=0$，所以由§1.4.2，
\[\boxed{C=\begin{bmatrix}2&1\\1&1\end{bmatrix}}.\]
$\det C=1>0$，向量非退化；$\operatorname{Cov}(Z_1,Z_2)=1$，两者不独立。
\key{按定义与独立性导出（§1.5.2）}
\[\begin{aligned}\phi_Z(u,v)&=E[e^{j(u(A+B)+vA)}]\\&=E[e^{j(u+v)A}]E[e^{juB}]\\&=\boxed{\exp[-\tfrac12((u+v)^2+u^2)]}.\end{aligned}\]
指数二次型等于$-\tfrac12(2u^2+2uv+v^2)$，交叉项系数不能漏掉2。
\key{三个核对}
$\phi_Z(0,0)=1$；$\phi_Z(u,0)=e^{-u^2}$对应方差2；$\phi_Z(0,v)=e^{-v^2/2}$对应方差1。联合特征函数不分解为两边缘乘积。
\end{tcolorbox}
\qtitle{第 27 题}{Day 9}
{\small\color{accentred}(Textbook 1.9 / Optional)}\par\smallskip
{\small\color{muted}教材小节：1.1.2}\par\medskip
\noindent\textbf{离散分布的冲激与右连续CDF}\par\smallskip\begin{tcolorbox}[breakable,colback=qbg,colframe=qbg,boxrule=0pt]\key{广义密度（§1.1.2）}
离散变量没有普通密度，教材所称密度可用冲激表示：
\[\boxed{f_X(x)=0.2\delta(x-1)+0.5\delta(x-2)+0.3\delta(x-3)}.\]
冲激的权重是概率，三者之和为1；不是有限高度的柱形密度。
\key{CDF逐点累积}
\[\boxed{F_X(x)=\begin{cases}0,&x<1,\\0.2,&1\le x<2,\\0.7,&2\le x<3,\\1,&x\ge3.\end{cases}}\]
CDF右连续，$F_X(2)=0.7$，不是0.2。
\key{图形}
\begin{center}\begin{tikzpicture}[x=1.2cm,y=2cm,>=stealth]
\draw[->](0,0)--(3.8,0)node[right]{$x$};\draw[->](0,0)--(0,1.3)node[above]{$F_X$};
\draw[thick](0,0)--(1,0);\draw[thick](1,.2)--(2,.2);\draw[thick](2,.7)--(3,.7);\draw[thick](3,1)--(3.7,1);
\foreach \x/\a/\b in {1/0/.2,2/.2/.7,3/.7/1}{\draw[fill=white](\x,\a)circle(1.7pt);\fill(\x,\b)circle(1.7pt);\node[below]at(\x,0){$\x$};}
\end{tikzpicture}\qquad\begin{tikzpicture}[x=1.2cm,y=3cm,>=stealth]
\draw[->](0,0)--(3.8,0)node[right]{$x$};\foreach\x/\p in{1/.2,2/.5,3/.3}{\draw[->,thick](\x,0)--(\x,\p)node[above]{$\p$};\node[below]at(\x,0){$\x$};}
\end{tikzpicture}\end{center}
\end{tcolorbox}
\qtitle{第 28 题}{Day 10}
{\small\color{accentred}(Textbook 1.11 / Optional)}\par\smallskip
{\small\color{muted}教材小节：1.1.2}\par\medskip
\noindent\textbf{Laplace密度的归一化与CDF}\par\smallskip\begin{tcolorbox}[breakable,colback=qbg,colframe=qbg,boxrule=0pt]\key{归一化}
\[1=a\int_{-\infty}^\infty e^{-|x|}dx=2a\int_0^\infty e^{-x}dx=2a,\quad\boxed{a=1/2}.\]
\key{分段积分（§1.1.2）}
$x<0$时积分区间内$|t|=-t$，所以$F(x)=\tfrac12\int_{-\infty}^xe^tdt=\tfrac12e^x$。
$x\ge0$时，先积负半轴再积$[0,x]$：
\[F(x)=\tfrac12+\tfrac12\int_0^xe^{-t}dt=1-\tfrac12e^{-x}.\]
\[\boxed{F(x)=\begin{cases}\tfrac12e^x,&x<0,\\1-\tfrac12e^{-x},&x\ge0.\end{cases}}\]
两段在0均为1/2，左端极限0、右端极限1，导数等于密度；不存在额外原点冲激。
\end{tcolorbox}
\qtitle{第 29 题}{Day 10}
{\small\color{accentred}(Textbook 1.28 / Optional)}\par\smallskip
{\small\color{muted}教材小节：1.5.1}\par\medskip
\noindent\textbf{独立和的特征函数}\par\smallskip\begin{tcolorbox}[breakable,colback=qbg,colframe=qbg,boxrule=0pt]\key{从定义证明相乘规则（§1.5.1）}
互相独立使复指数的期望因子分解：
\[\phi_{b+\sum a_kX_k}(v)=E[e^{jvb}\prod e^{jva_kX_k}]=e^{jvb}\prod_k\phi_k(a_kv).\]
零系数的因子是$\phi_k(0)=1$。四问依次为
\[\boxed{\phi_1(v)\phi_2(v)},\quad\boxed{\phi_1(v)\phi_2(v)\phi_3(v)},\]
\[\boxed{\phi_1(v)\phi_2(2v)\phi_3(3v)},\quad\boxed{e^{j10v}\phi_1(2v)\phi_2(v)\phi_3(4v)}.\]
缩放改变自变量，平移产生相位因子；二者都不是直接乘以随机变量的系数。若只有两两独立，三因子一般不能如此分解。
\end{tcolorbox}
\qtitle{第 30 题}{Day 10}
{\small\color{accentred}(Textbook 1.30 / Optional)}\par\smallskip
{\small\color{muted}教材小节：1.5.1}\par\medskip
\noindent\textbf{均匀分布的Fourier与特征函数约定}\par\smallskip\begin{tcolorbox}[breakable,colback=qbg,colframe=qbg,boxrule=0pt]\key{先写密度和符号约定}
$X\sim U(a,b)$，$a<b$，$f_X=1/(b-a)$在$[a,b]$上。
若$\mathcal F\{f\}(\omega)=\int f(x)e^{-j\omega x}dx$，则由§1.5.1，$\phi_X(v)=\mathcal F\{f_X\}(-v)$，不是同号代入。
\key{直接积分}
\[\begin{aligned}\phi_X(v)&=\frac1{b-a}\int_a^b e^{jvx}dx\\&=\frac{e^{jvb}-e^{jva}}{jv(b-a)}\\&=\boxed{e^{jv(a+b)/2}\frac{\sin(v(b-a)/2)}{v(b-a)/2}},\quad v\ne0.\end{aligned}\]
在$v=0$由定义$\phi_X(0)=1$，也等于上式的连续极限。
\key{常见特例}
$U(-A,A)$得到$\sin(Av)/(Av)$；$U(0,A)$得到$e^{jAv/2}\sin(Av/2)/(Av/2)$。相位因子编码均值，不能随意删去。
\end{tcolorbox}
\end{document}